% less than one page
\section{Related Work}~\label{sec:related}
% This is common among out-of-place systems, including LSM-tree engines 
% (RocksDB~\cite{rocksdb-saf}, LevelDB~\cite{leveldb}, Pebble~\cite{pebble}, TigerBeetle~\cite{tigerbeetle-arch}), 
% log-structured file systems (F2FS~\cite{DBLP:conf/fast/LeeSHC15}, NILFS2~\cite{nilfs2}, LFS~\cite{lfs}), 
% zoned storage engines for ZNS SSDs or host-managed SMR (ZenFS~\cite{zenfs} or SMRSTORE~\cite{DBLP:conf/usenix/Hwang0KES25}), 
% and cloud storage backends (Aurora~\cite{aurora}, S3 shardstore~\cite{shardstore}, Ceph BlueStore~\cite{bluestore}).

\module{Different ways to perform out-of-place writes.}
\revision[R3.W2]{ Out-of-place writes appear in several forms.
WiredTiger~\cite{wiredtiger-compaction} stores each table in a cyclic file and rewrites valid pages to the file head during compaction.
LSM engines~\cite{rocksdb-saf, leveldb, tigerbeetle-arch} also write out-of-place, reclaiming space only when compaction is triggered, making GC timing and reclaimed space hard to control.
In contrast, ZLeanStore reclaims space proactively at zone granularity~\cite{DBLP:conf/fast/LeeSHC15, nilfs2, lfs, zenfs, DBLP:conf/usenix/Hwang0KES25, skylight}, independent of other DBMS components~\cite{DBLP:journals/tos/LuPGAA17}.
Doing so at the DBMS layer allows GC to account for both DB- and SSD-level WA. }

\module{Optimizing data structures to reduce DB WA.}
A substantial amount of work has reduced WA by optimizing index structures.
B-tree-based storage engines such as WO-B-tree~\cite{WObtree}, Bf-tree~\cite{DBLP:journals/pvldb/HaoC24}, and B-epsilon tree~\cite{DBLP:journals/usenix-login/BenderFJJKPY015} buffer updates and apply them in batches to reduce index-induced WA.
LSM-tree engines were introduced to address this, and various works attempt to optimize it further~\cite{rocksdb-saf, leveldb, tigerbeetle_github, novelsm,DBLP:conf/edbt/VinconHRO0P18,LSMTreeCost_SIGMOD2024,Dayan_TODS2018,Liu2025_ArceKV,lsmdesign,Dayan_VLDB2022_Hybrid}.
These approaches are orthogonal to ours: they reduce WA at the index level, whereas our techniques reduce DB WA at the page level without changing the index structure.
Thus, they can be complementary when applied together.

\module{Reducing WA by mitigating space amplification.}
Write volume can be reduced by lowering overall storage consumption.
Techniques such as XMerge~\cite{DBLP:conf/cidr/AlhomssiL21} and B-tree node compaction~\cite{DBLP:journals/access/LeeAL23} 
improve node utilization and reduce disk writes; these index-level optimizations are orthogonal and can benefit from our approach.
\revision[R2.D4]{Disk-level compression in MySQL~\cite{mysql-comp}, WiredTiger~\cite{wiredtiger-compaction}, and RocksDB~\cite{rocksdb-saf} 
similarly reduces write volume and storage footprint, 
but is ineffective under 4~KiB page granularity, whereas we provide a solution for 4~KiB pages. }
Moreover, out-of-place--write engines do not examine how compression interacts with DB WAF induced by GC, whereas we explicitly analyze and address these effects.
Some SSDs also employ internal compression to increase overprovisioning and mitigate WA~\cite{DBLP:conf/osdi/ZuckTSH14}; 
our DB GC naturally benefits from compression as well, but, unlike SSDs, we can flexibly enable, disable, or tune compression based on application needs.


\module{Mitigating WA upon GC.}
Prior work proposes a wide range of GC algorithms~\cite{DBLP:conf/eurosys/ZhouWHHXZ15, DBLP:journals/corr/DayanB15, DBLP:conf/fast/MinKCLE12, DBLP:journals/cacm/LangeNY23, DBLP:journals/tos/Desnoyers14, DBLP:conf/fast/LeeSHC15, lockv, DBLP:journals/tos/LuPGAA17, DBLP:journals/pvldb/YuMKLMK22, DBLP:journals/pe/Houdt14, DBLP:conf/hotstorage/ShafaeiDF16} 
that classify pages by lifetime, recency, or spatial locality to reduce WA during GC.
Our approach differs in that we predict deathtimes of individual pages using database-specific semantics, 
which lower layers such as SSD GC cannot exploit~\cite{DBLP:journals/cacm/LangeNY23,DBLP:conf/eurosys/HeKAA17, DBLP:journals/pomacs/LangeNY25}.
To our knowledge, this is the first work to integrate GDT-based GC within a DBMS, leveraging workload knowledge that is only available to the DBMS.

\module{Optimizing system write patterns for SSDs.}
Prior work has examined how applications can tailor I/O to SSD characteristics to exploit internal parallelism~\cite{DBLP:conf/eurosys/HeKAA17, DBLP:journals/pvldb/KakaraparthyPPK19, DBLP:journals/pvldb/HeXLJZYML23}, analyzed tree-based access patterns~\cite{DBLP:journals/pvldb/DidonaISK20}, and characterized SSD behavior under diverse workloads~\cite{DBLP:journals/tos/YadgarGJS21, DBLP:conf/hotstorage/YadgarG16}.
Although these studies improve latency or throughput, they do not explicitly target SSD WAF; even when WAF is considered, it is not fully eliminated~\cite{DBLP:conf/osdi/YangPGTS14, DBLP:conf/fast/Curtis-MauryKRM24}.
In contrast, our DBMS GC directly minimizes SSD WAF while accounting for its interaction with DB WAF to reduce total WAF.

\module{Zoned storage for ZNS.}
Several systems~\cite{ssdfs, zenfs, DBLP:conf/usenix/Hwang0KES25, yang2024zonedstorageoptimizedflash, DBLP:conf/hotstorage/LeeLL022, DBLP:journals/fgcs/HaS25} 
leverage ZNS SSDs~\cite{DBLP:conf/usenix/BjorlingAHRMGA21} to exploit the SSD WAF of~1.
While effective, these designs are tied to ZNS hardware and thus limited to environments where such devices are available.
In contrast, our approach treats ZNS as an optional optimization rather than a requirement, making it applicable across a broader range of devices.

\module{Systems using FDP.}
The FDP interface was introduced only recently~\cite{nvme_fdp_tp4146}, and emerging work explores its potential for reducing WAF~\cite{DBLP:conf/hotstorage/ZhangPCK23,eurosys25,george2024fdp}.
Our work identifies the optimal use case for FDP placement hints from the perspective of out-of-place DBMSs to achieve an SSD WAF of~1.
We further suggest appropriate settings for DB GC granularity and the number of active zones under FDP.
